\documentclass[11pt,a4paper]{article}

\usepackage[margin=2.5cm]{geometry}
\usepackage{graphicx}
\usepackage{booktabs}
\usepackage{array}
\usepackage{longtable}
\usepackage{hyperref}
\usepackage{xcolor}
\usepackage{float}
\usepackage{listings}
\usepackage{caption}
\usepackage{amssymb}

\hypersetup{
  colorlinks=true,
  linkcolor=blue,
  urlcolor=blue,
  citecolor=blue
}

\graphicspath{{screenshots/}{portfolio/screenshots/}}

\lstdefinestyle{code}{
  basicstyle=\ttfamily\small,
  breaklines=true,
  frame=single,
  columns=fullflexible
}
\lstset{style=code}

\newcommand{\projecttitle}{Analysing the AI Investment Landscape using Knowledge Graphs}
\newcommand{\studentname}{Michael Pertl}
\newcommand{\coursename}{Knowledge Graphs}
\newcommand{\submissiondate}{\today}
\newcommand{\checkedbox}{\(\boxtimes\)}
\newcommand{\emptybox}{\(\square\)}

\newcommand{\screenshot}[2]{
  \begin{figure}[H]
    \centering
    \IfFileExists{#1}{
      \includegraphics[width=0.92\textwidth]{#1}
    }{
      \fbox{
        \begin{minipage}[c][6cm][c]{0.86\textwidth}
          \centering
          Screenshot placeholder\\[0.4em]
          \texttt{#1}\\[0.4em]
          #2
        \end{minipage}
      }
    }
    \caption{#2}
  \end{figure}
}

\begin{document}

\begin{titlepage}
  \centering
  \vspace*{2cm}
  {\LARGE \textbf{Knowledge Graphs Mini-Project Portfolio}\par}
  \vspace{1.5cm}
  {\Large \projecttitle\par}
  \vspace{2cm}
  \begin{tabular}{rl}
    Project Title: & \projecttitle \\
    Name: & \studentname \\
    Course: & \coursename \\
    Mode: & \checkedbox{} 6 ECTS \quad \emptybox{} 3 ECTS \\
    Date: & \submissiondate \\
  \end{tabular}
  \vfill
  \begin{minipage}{0.86\textwidth}
    \small
    This document follows the portfolio pro forma: first a learning-outcome
    self-assessment and declaration, then the actual mini-project report.
    Replace all placeholder text before submission.
  \end{minipage}
\end{titlepage}

\section*{Learning Outcome Self-Assessment}
\addcontentsline{toc}{section}{Learning Outcome Self-Assessment}

The project mainly addresses real-world and financial Knowledge Graph
applications. The primary focus is on LO9 and LO10. The portfolio also shows
basic proficiency in LO4, LO5, LO6, LO7, LO11, and LO12. LO1 and LO3 are not
claimed because the project does not implement Knowledge Graph Embeddings or
Graph Neural Networks. LO2 and LO8 are also not claimed as separate learning
outcomes because logical knowledge representation and KG evolution are not the
central contribution of this mini-project.

\subsection*{Representations}

\begin{longtable}{p{0.11\textwidth}p{0.39\textwidth}p{0.16\textwidth}p{0.24\textwidth}}
\toprule
LO & Learning outcome & Claim & Where discussed \\
\midrule
\endhead
LO1 & Understand and apply Knowledge Graph Embeddings &
\emptybox{} basic \quad \emptybox{} exceeded &
Not claimed. The project deliberately focuses on property graph modelling,
Cypher analysis, and application design rather than KGE models such as TransE,
ComplEx, or ConvE. \\
\addlinespace
LO2 & Understand and apply logical knowledge in KGs &
\emptybox{} basic \quad \emptybox{} exceeded &
Not claimed as a separate LO. The project uses constraints and validation
rules, but it does not implement recursive logical reasoning or existential
object creation as a main method. \\
\addlinespace
LO3 & Understand and apply Graph Neural Networks &
\emptybox{} basic \quad \emptybox{} exceeded &
Not claimed. The project does not train or apply a GNN; graph structure is used
through Neo4j queries and network analysis rather than neural message passing. \\
\addlinespace
LO4 & Compare different Knowledge Graph data models from database, semantic web, machine learning and data science communities &
\checkedbox{} basic \quad \emptybox{} exceeded &
Basic proficiency: the portfolio explains why a Neo4j property graph was used
instead of a purely tabular model, and contrasts this choice with RDF/Semantic
Web and ML-oriented graph representations in Section~\ref{sec:model}. \\
\bottomrule
\end{longtable}

\subsection*{Systems}

\begin{longtable}{p{0.11\textwidth}p{0.39\textwidth}p{0.16\textwidth}p{0.24\textwidth}}
\toprule
LO & Learning outcome & Claim & Where discussed \\
\midrule
\endhead
LO5 & Design and implement architectures of a Knowledge Graph &
\checkedbox{} basic \quad \emptybox{} exceeded &
Basic proficiency: the project implements a KG architecture from raw evidence
CSV files to processed tables, Neo4j import files, Cypher constraints, analysis
queries, and an optional web application. See Sections~\ref{sec:model},
\ref{sec:pipeline}, and~\ref{sec:neo4j}. \\
\addlinespace
LO6 & Describe and apply scalable reasoning methods in Knowledge Graphs &
\checkedbox{} basic \quad \emptybox{} exceeded &
Basic proficiency: the project applies scalable graph querying in Neo4j to
derive regional funding totals, domain concentration, investor participation,
co-investment pairs, and cross-region investor links. See Sections~\ref{sec:neo4j}
and~\ref{sec:results}. \\
\addlinespace
LO7 & Apply a system to create a Knowledge Graph &
\checkedbox{} basic \quad \emptybox{} exceeded &
Basic proficiency: the project creates a Neo4j KG from curated public funding
evidence using Python normalization scripts and Cypher import queries. See
Sections~\ref{sec:pipeline} and~\ref{sec:neo4j}. \\
\addlinespace
LO8 & Apply a system to evolve a Knowledge Graph &
\emptybox{} basic \quad \emptybox{} exceeded &
Not claimed as a separate LO. The repository supports refreshable data
generation, but the portfolio does not focus on KG completion, rule learning,
link prediction, or schema evolution. \\
\bottomrule
\end{longtable}

\subsection*{Applications}

\begin{longtable}{p{0.11\textwidth}p{0.39\textwidth}p{0.16\textwidth}p{0.24\textwidth}}
\toprule
LO & Learning outcome & Claim & Where discussed \\
\midrule
\endhead
LO9 & Describe and design real-world applications of Knowledge Graphs &
\checkedbox{} basic \quad \emptybox{} exceeded &
Main focus: the project designs a real-world KG application for analysing the
AI investment ecosystem and provides an optional graph explorer. See
Sections~\ref{sec:introduction},~\ref{sec:results}, and~\ref{sec:webapp}. \\
\addlinespace
LO10 & Describe financial Knowledge Graph applications &
\checkedbox{} basic \quad \emptybox{} exceeded &
Main focus: the project applies KG methods to venture funding and investment
network analysis, including investor concentration, co-investment structure,
regional comparison, and sector funding flows. See Sections~\ref{sec:scope}
and~\ref{sec:results}. \\
\addlinespace
LO11 & Apply a system to provide services through a Knowledge Graph &
\checkedbox{} basic \quad \emptybox{} exceeded &
Basic proficiency: the React web application and graph export show how the KG
can provide an exploration service beyond direct Cypher queries. See
Section~\ref{sec:webapp}. \\
\addlinespace
LO12 & Describe connections between KGs, ML and AI &
\checkedbox{} basic \quad \emptybox{} exceeded &
Basic proficiency: the project studies AI companies using KG techniques and
discusses how symbolic graph analysis differs from, and could complement,
ML-based approaches such as embeddings or GNNs. See Sections~\ref{sec:introduction}
and~\ref{sec:conclusion}. \\
\bottomrule
\end{longtable}

\newpage
\section*{Additional Information}
\addcontentsline{toc}{section}{Additional Information}

This information has no effect on marking according to the pro forma. Fill it
out honestly before submission.

\begin{table}[H]
\centering
\begin{tabular}{p{0.62\textwidth}p{0.25\textwidth}}
\toprule
Question & Answer \\
\midrule
How many hours did you spend on your mini-project? & \texttt{<XX>} hours \\
How many hours did you spend on portfolio document preparation? & \texttt{<XX>} hours \\
\bottomrule
\end{tabular}
\caption{Time-spent declaration placeholders.}
\end{table}

\section*{Declaration}
\addcontentsline{toc}{section}{Declaration}

\emptybox{} I confirm that I have marked all parts generated by Generative AI
or otherwise copied from other sources and given any prompt I used either in a
footnote or in an appendix, making clear which parts are generated by which
prompts or similar.

\vspace{1em}
\noindent
Generative AI use in the mini-project:
\emptybox{} I used it for roughly \texttt{<XX>\%} of the mini-project.

\vspace{0.5em}
\noindent
\textit{Description:} \texttt{<Describe which parts were generated with AI and
which parts you implemented, checked, curated, interpreted, or wrote yourself.>}

\vspace{1em}
\noindent
Generative AI use in this document:
\emptybox{} I used it for roughly \texttt{<XX>\%} of this document.

\vspace{0.5em}
\noindent
\textit{Description:} \texttt{<Describe AI assistance for drafting, structure,
tables, wording, or LaTeX formatting.>}

\newpage
\tableofcontents
\newpage

\section*{Portfolio Report}
\addcontentsline{toc}{section}{Portfolio Report}

In paragraphs where a learning outcome is explicitly discussed, references such
as (LO5), (LO7), or (LO10) are included to match the pro forma instructions.

\section{Introduction and Motivation}
\label{sec:introduction}

\subsection{Motivation}

The AI ecosystem is growing rapidly, with complex relationships between
startups, venture capital firms, corporate investors, countries, regions,
funding rounds, and technological domains. These relationships are difficult to
analyse using traditional tabular data alone because many of the relevant
questions are relational: which investors repeatedly co-invest, which domains
attract the largest funding rounds, which companies connect multiple investor
groups, and how regional funding patterns develop over time.

\subsection{Problem}

The problem addressed in this mini-project is that public AI funding
information is fragmented across news articles, press releases, trackers, and
company announcements. This makes it challenging to identify investor focus,
venture-capital collaboration patterns, regional differences, and emerging AI
application areas based on where large amounts of capital flow.

\subsection{Solution}

The solution is a self-curated Knowledge Graph application for analysing AI
mega-round investment networks (LO9, LO10). The graph integrates public funding
events, investors, AI startups, AI domains, countries, regions, and time using a
Neo4j property graph. This allows structured Cypher queries over investor
participation, sector concentration, co-investment relationships, temporal
development, and regional comparison.

\subsection{Research Question}

\begin{quote}
How do major regional AI mega-round funding networks differ in terms of investor
concentration, sector focus, co-investment structure, and temporal development
from 2024 to 2026?
\end{quote}

\subsection{Project Goal}

The project builds a reproducible data pipeline and Neo4j property graph for
analysing public AI mega-round funding events. It also includes an optional web
application as an exploration layer, but the analytical source of truth is the
curated CSV data and Neo4j graph.

\subsection{Implemented Steps}

The project follows these practical steps:

\begin{enumerate}
  \item collect and curate public funding evidence;
  \item filter for AI-native or AI-infrastructure companies and USD 100M+
        mega-rounds;
  \item classify startups into AI domains and assign countries and regions;
  \item build a Neo4j property graph with startups, investors, funding rounds,
        domains, countries, and regions;
  \item analyse sector concentration, investor networks, co-investment pairs,
        regional differences, and temporal development;
  \item provide an optional web application as a graph exploration service.
\end{enumerate}

\section{Dataset Scope and Limitations}
\label{sec:scope}

\subsection{Inclusion Criteria}

A funding event is included only if all of the following conditions are met:

\begin{enumerate}
  \item The company is AI-native or AI-infrastructure-related.
  \item The company headquarters or primary operating identity can be assigned
        to a documented country and region.
  \item The funding event occurred between 2024 and 2026.
  \item The round amount is at least USD 100 million or equivalent.
  \item A public source confirms the amount, date or year, and at least one
        investor.
  \item The event has a source URL.
\end{enumerate}

\subsection{Current Dataset Size}

\begin{table}[H]
\centering
\begin{tabular}{lr}
\toprule
Entity or relationship & Count \\
\midrule
Funding events & 201 \\
Startups & 160 \\
Investors & 520 \\
Investor participations & 1,156 \\
AI domains & 14 \\
Countries & 15 \\
Regions & 9 \\
\bottomrule
\end{tabular}
\caption{Current validated dataset size.}
\end{table}

\subsection{Regional Coverage}

\begin{table}[H]
\centering
\begin{tabular}{lrrr}
\toprule
Region & 2024 rounds & 2025 rounds & 2026 rounds \\
\midrule
US & 52 & 63 & 35 \\
Europe & 8 & 15 & 9 \\
China & 4 & 0 & 1 \\
Israel & 1 & 5 & 1 \\
Canada & 2 & 1 & 0 \\
UAE / Gulf & 1 & 0 & 0 \\
India & 0 & 1 & 0 \\
Singapore & 0 & 0 & 1 \\
Other Global & 1 & 0 & 0 \\
\bottomrule
\end{tabular}
\caption{Funding event counts by region and year.}
\end{table}

\subsection{Important Caveat}

This graph represents curated publicly reported AI mega-rounds above the USD
100 million threshold. It does not represent the full AI startup ecosystem. It
is biased toward visible late-stage and capital-intensive companies. Therefore,
the findings should be interpreted as statements about this curated public
mega-round dataset, not as claims about all AI investment.

\section{Knowledge Graph Model}
\label{sec:model}

\subsection{Core Model}

The graph uses funding rounds as first-class nodes because each round has its
own amount, date, round type, source URL, confidence level, and investor
participation list.

The chosen representation is a Neo4j property graph (LO4). This is suitable for
the project because funding events and investor participations naturally have
properties such as amount, date, role, confidence, and source URL. A purely
tabular model would make co-investment and neighbourhood queries less direct.
An RDF/Semantic Web model could also represent the domain, but the mini-project
focuses on exploratory graph analytics with Cypher rather than ontology
reasoning. Machine-learning graph representations such as embeddings or GNNs
are not used because the project scope emphasizes transparent data modelling
and interpretable financial network analysis.

\begin{lstlisting}[language={}]
(:Startup)-[:RAISED]->(:FundingRound)
(:Investor)-[:PARTICIPATED_IN {role}]->(:FundingRound)
(:Startup)-[:OPERATES_IN {role}]->(:AIDomain)
(:Startup)-[:HEADQUARTERED_IN]->(:Country)
(:Country)-[:PART_OF]->(:Region)
\end{lstlisting}

\subsection{Node Types}

\begin{table}[H]
\centering
\begin{tabular}{ll}
\toprule
Node label & Purpose \\
\midrule
Startup & AI startup or AI infrastructure company \\
FundingRound & Publicly reported qualifying funding event \\
Investor & VC, corporate investor, investment manager, or other participant \\
AIDomain & Primary or secondary AI sector classification \\
Country & Documented company country \\
Region & Analytical regional grouping \\
\bottomrule
\end{tabular}
\caption{Main graph node types.}
\end{table}

\subsection{Relationship Types}

\begin{table}[H]
\centering
\begin{tabular}{ll}
\toprule
Relationship & Meaning \\
\midrule
RAISED & Startup raised a funding round \\
PARTICIPATED\_IN & Investor participated in a funding round \\
OPERATES\_IN & Startup belongs to an AI domain \\
HEADQUARTERED\_IN & Startup is assigned to a country \\
PART\_OF & Country belongs to an analytical region \\
\bottomrule
\end{tabular}
\caption{Main graph relationship types.}
\end{table}

\section{Data Pipeline}
\label{sec:pipeline}

\subsection{Pipeline Overview}

The pipeline is deterministic and CSV-based. Raw evidence is kept in
\texttt{data/raw/}. Processed normalized tables are written to
\texttt{data/processed/}. Neo4j import-ready CSV files are written to
\texttt{data/neo4j\_import/}.

This architecture separates evidence collection, validation, normalization,
graph import, and application export (LO5, LO7). Raw data remains auditable,
processed data is reproducible, and the Neo4j graph can be rebuilt from the
CSV pipeline.

\begin{lstlisting}[language={}]
data/raw/funding_events_raw.csv
        |
        v
scripts/validate_raw_data.py
        |
        v
scripts/normalize_raw_data.py
        |
        v
data/processed/*.csv
        |
        v
data/neo4j_import/*.csv
        |
        v
Neo4j import using queries/00_schema.cypher and queries/01_import.cypher
\end{lstlisting}

\subsection{Main Commands}

\begin{lstlisting}[language=bash]
python3 scripts/validate_raw_data.py
python3 scripts/normalize_raw_data.py
python3 scripts/build_growth_metrics.py
python3 scripts/report_data_quality.py
python3 -m unittest discover -s tests -v
\end{lstlisting}

\subsection{Validation Rules}

The validation script checks required fields, controlled region vocabulary,
valid year range, minimum funding threshold, source URLs, investor presence,
domain presence, and duplicate event identifiers. Rows are not silently dropped.

\section{Neo4j Import and Validation}
\label{sec:neo4j}

\subsection{Starting Neo4j}

\begin{lstlisting}[language=bash]
docker compose up -d
\end{lstlisting}

Neo4j Browser:

\begin{lstlisting}[language={}]
http://localhost:7474
username: neo4j
password: startupkg2026
\end{lstlisting}

\subsection{Import Commands}

\begin{lstlisting}[language=bash]
docker exec -i startupkg-neo4j cypher-shell -u neo4j -p startupkg2026 < queries/00_schema.cypher
docker exec -i startupkg-neo4j cypher-shell -u neo4j -p startupkg2026 < queries/01_import.cypher
docker exec -i startupkg-neo4j cypher-shell -u neo4j -p startupkg2026 < queries/02_validation.cypher
\end{lstlisting}

\subsection{Validation Screenshot}

\screenshot{neo4j_validation_counts.png}{Neo4j validation counts for nodes and relationships.}

\section{Analysis Results}
\label{sec:results}

This section should use result tables from \texttt{queries/04\_final\_results.cypher}.
Replace or extend the tables below with exported Neo4j Browser results.

\subsection{Querying and Derived Insights}

The project applies Neo4j and Cypher as the main reasoning and analytics system
(LO6). The queries do not infer hidden facts with embeddings or GNNs. Instead,
they derive interpretable graph-level insights from explicit evidence:
aggregated funding volume, repeated investor participation, co-investment
pairs, cross-region investor activity, and startup funding trajectories. This
matches the project scope because the key requirement is transparent analysis
of a curated financial evidence graph.

\subsection{Funding by Region and Year}

\begin{table}[H]
\centering
\begin{tabular}{llrr}
\toprule
Region & Year & Rounds & Amount USD bn \\
\midrule
US & 2024 & 52 & 43.16 \\
US & 2025 & 63 & 78.20 \\
US & 2026 & 35 & 237.54 \\
Europe & 2024 & 8 & 4.00 \\
Europe & 2025 & 15 & 6.36 \\
Europe & 2026 & 9 & 6.68 \\
China & 2024 & 4 & 2.71 \\
China & 2026 & 1 & 2.00 \\
\bottomrule
\end{tabular}
\caption{Selected funding volume by region and year.}
\end{table}

\subsection{Domain Totals}

\begin{table}[H]
\centering
\begin{tabular}{lrrr}
\toprule
AI domain & Startups & Rounds & Amount USD bn \\
\midrule
Foundation Models / LLMs & 25 & 36 & 319.39 \\
AI Infrastructure & 27 & 32 & 22.24 \\
Enterprise AI & 22 & 28 & 6.50 \\
Coding Agents / Developer Tools & 9 & 11 & 6.21 \\
Robotics / Physical AI & 11 & 13 & 5.84 \\
AI Chips / Semiconductors & 15 & 17 & 5.74 \\
\bottomrule
\end{tabular}
\caption{Largest AI domains by funding volume in the curated dataset.}
\end{table}

\subsection{Most Active Investors}

\begin{table}[H]
\centering
\begin{tabular}{lrr}
\toprule
Investor & Participated rounds & Lead rounds \\
\midrule
NVIDIA & 39 & 3 \\
Lightspeed Venture Partners & 29 & 8 \\
Andreessen Horowitz & 27 & 10 \\
Sequoia Capital & 25 & 8 \\
General Catalyst & 24 & 8 \\
\bottomrule
\end{tabular}
\caption{Selected most active investors by participation count.}
\end{table}

\subsection{Preliminary Interpretation}

Within this curated public mega-round dataset, funding volume is heavily
concentrated in US-based companies and in foundation model companies. Investor
participation also appears concentrated around a small group of recurring VC,
corporate, and investment-manager actors. Cross-region participation suggests
that several US and global investors connect otherwise regional funding
subgraphs (LO9, LO10).

These findings must be interpreted carefully because 2026 is partial-year data
and because several large strategic funding rounds strongly affect total funding
volume.

\section{Graph Visualizations}
\label{sec:visualizations}

Use \texttt{queries/05\_visualization\_views.cypher} and switch the Neo4j
Browser result view to Graph. Add screenshots below.

\subsection{AI Infrastructure Network}

\screenshot{neo4j_ai_infrastructure_network.png}{AI Infrastructure startup, round, and investor network.}

\subsection{Foundation Models / LLM Network}

\screenshot{neo4j_foundation_models_network.png}{Foundation Models / LLM startup, round, and investor network.}

\subsection{Cross-Region Investor Network}

\screenshot{neo4j_cross_region_investors.png}{Investors participating across multiple startup regions.}

\subsection{Startup Neighborhood}

\screenshot{neo4j_mistral_neighborhood.png}{Mistral AI neighborhood with rounds, investors, domains, country, and region.}

\section{Optional Web Application}
\label{sec:webapp}

In addition to the Neo4j analysis workflow, the project includes a React-based
web application for exploring the graph export. The web app is an optional
interface for inspecting companies, investors, regions, domains, and market
events. This demonstrates how a KG can provide a service to users through
visual exploration and filtering (LO11). It is not a replacement for the Neo4j
analytical layer; it is an additional presentation and exploration interface.

\subsection{Running the Web App}

\begin{lstlisting}[language=bash]
./start.sh
\end{lstlisting}

or manually:

\begin{lstlisting}[language=bash]
cd web
npm install
npm run dev
\end{lstlisting}

\subsection{Web App Screenshots}

\screenshot{web_app_graph_explorer.png}{React web app graph explorer.}

\screenshot{web_app_company_detail.png}{React web app company or filtered detail view.}

\section{Conclusion}
\label{sec:conclusion}

The project shows how a knowledge graph can support analysis of public AI
mega-round funding networks. The graph connects startups, funding rounds,
investors, AI domains, countries, regions, and time. This structure supports
regional comparison, investor concentration analysis, co-investment analysis,
domain-level funding comparison, and company-level funding trajectories.

The main limitation is that the dataset is curated and threshold-based. It
captures large publicly reported AI funding events, not the full AI startup or
AI investment landscape.

The project also illustrates the connection between Knowledge Graphs, AI, and
Machine Learning (LO12). The subject matter is the AI investment ecosystem, but
the method is symbolic and graph-based rather than ML-based. This makes the
results explainable: every relationship can be traced back to a public funding
event and source URL. Future work could add ML-oriented KG methods such as
link prediction or embeddings, but those are intentionally outside the current
scope.

\section{Appendix}

\subsection{Repository Structure}

\begin{lstlisting}[language={}]
docs/
data/raw/
data/processed/
data/neo4j_import/
scripts/
queries/
tests/
web/
app/
\end{lstlisting}

\subsection{Important Files}

\begin{itemize}
  \item \texttt{data/raw/funding\_events\_raw.csv}: core evidence table.
  \item \texttt{scripts/validate\_raw\_data.py}: raw data validation.
  \item \texttt{scripts/normalize\_raw\_data.py}: normalized CSV generation.
  \item \texttt{queries/00\_schema.cypher}: constraints and indexes.
  \item \texttt{queries/01\_import.cypher}: Neo4j CSV import.
  \item \texttt{queries/04\_final\_results.cypher}: final analysis tables.
  \item \texttt{queries/05\_visualization\_views.cypher}: graph screenshots.
\end{itemize}

\subsection{Suggested Screenshot Filenames}

Place screenshots in \texttt{portfolio/screenshots/} with these names:

\begin{lstlisting}[language={}]
neo4j_validation_counts.png
neo4j_ai_infrastructure_network.png
neo4j_foundation_models_network.png
neo4j_cross_region_investors.png
neo4j_mistral_neighborhood.png
web_app_graph_explorer.png
web_app_company_detail.png
\end{lstlisting}

\subsection{Careful Wording}

Use cautious wording in the final report:

\begin{quote}
Within this curated public AI mega-round dataset, the graph suggests ...
\end{quote}

Avoid wording that implies complete market coverage.

\end{document}
